%% kapa: exam=1 round=30 year=2019 period=1 form=A subject=civil_law subject_ko=민법 seq=20 number=20
%% source: Q-Net 감정평가사 기출문제 (cst00309) artlSeq=5205205 file=extracted/1008/5205205/2019년 제30회 감정평가사 1차시험 문제지/2019년 제30회 감정평가사 1차시험 문제지/2019년도 제30회 감정평가사 1차 1교시 A형.hwp
%% conversion: kordoc 4.24.51 -> tools/gen1.py (auto); review=proofread-sonnet; flags=none
\begin{question}{20}
소멸시효의 기산점에 관한 설명으로 옳지 않은 것은? (다툼이 있으면 판례에 따름)
\begin{choices}
  \item 소멸시효는 권리를 행사할 수 있는 때로부터 진행하며, 이때 ‘권리를 행사할 수 있다’는 것은 권리를 행사함에 있어 원칙적으로 법률상 장애가 없는 것을 가리킨다.
  \item 부작위를 목적으로 하는 채권의 소멸시효는 위반행위를 한 때로부터 진행한다.
  \item 정지조건부권리의 경우에는 조건 미성취의 동안은 권리를 행사할 수 없는 것이어서 소멸시효가 진행되지 않는다.
  \item 소유권이전등기의무의 이행불능으로 인한 전보배상청구권의 소멸시효는 이전등기의무가 이행불능이 된 때부터 진행된다.
  \item 본래의 소멸시효 기산일과 당사자가 주장하는 기산일이 서로 다른 경우에는 법원은 본래의 소멸시효 기산일을 기준으로 소멸시효를 계산하여야 한다.
\end{choices}
\end{question}
\answer{5}
